In Kazakhstan and similar data-scarce urban settings, official city totals may be available while true cell-level census labels are not. This creates a two-part methodological problem: first, how to construct a reproducible calibrated proxy population surface under missing cell-level truth; second, how to prevent that deterministic surface from being over-read as equally reliable everywhere. City1 addresses both stages within one framework. Stage 1 builds a fixed 500 m calibrated proxy population surface from OpenStreetMap-derived features, weak target construction, Random Forest modeling, and official-total calibration. Stage 2 adds a reliability-aware interpretation layer through calibrated ensemble predictions, P10/P50/P90 proxy intervals, uncertainty width, relative uncertainty, confidence bands, and hotspot prioritization classes. The baseline stage is supported by leave-one-city-out transfer testing, spatial block cross-validation, grid benchmarking, OSM completeness assessment, district aggregation evidence, external structural comparison, ablation, and qualitative plausibility review. The uncertainty-aware stage is supported most strongly by hotspot stability and confidence-band separation of uncertainty burden; interval coverage is mixed, external disagreement alignment is mixed or negative, and district interval coverage remains partial. The framework therefore does not reconstruct true cell-level census counts and does not estimate true census uncertainty. Instead, it contributes a transparent, reproducible, officially calibrated, and uncertainty-aware proxy population framework for screening and interpretation under missing cell-level truth.
